% 图契约：源尺寸、观察距离和电流机制分别由三个无量纲量判断；几何、方向和公式与正文同一约定。
\documentclass[tikz,border=5pt]{standalone}
\usepackage{fontspec,xeCJK,amsmath,amssymb}
\IfFontExistsTF{Times New Roman}{\setmainfont{Times New Roman}}{\setmainfont{TeX Gyre Termes}}
\IfFontExistsTF{Heiti SC}{\setCJKmainfont[AutoFakeBold=1.5]{Heiti SC}}{\setCJKmainfont[AutoFakeBold=1.5]{FandolHei-Regular}}
\usetikzlibrary{arrows.meta,calc,patterns}
\definecolor{ink}{HTML}{243745}
\definecolor{blue}{HTML}{275DAB}
\definecolor{green}{HTML}{177C63}
\definecolor{red}{HTML}{BD4935}
\definecolor{orange}{HTML}{AA6C16}
\definecolor{purple}{HTML}{715393}
\begin{document}
\begin{tikzpicture}[x=1cm,y=1cm,font=\fontsize{9.5}{12}\selectfont,
v/.style={-{Latex[length=2mm,width=1.3mm]},line width=.85pt},
axis/.style={-{Latex[length=1.7mm,width=1.1mm]},ink!65,line width=.55pt},
aux/.style={ink!45,dashed,line width=.55pt},
lab/.style={fill=white,inner sep=1.8pt},
card/.style={draw=ink!25,fill=blue!3,rounded corners=2pt,text width=6.1cm,minimum height=.85cm,inner sep=6pt,align=center}]
\path[use as bounding box](-8,-3.6)rectangle(8,3.65);
\node[font=\bfseries\fontsize{10.5}{13}\selectfont,text width=15.1cm,align=center]at(0.00000,3.22000){M7-H47　源尺寸、观察距离和电流机制分别由三个无量纲量判断};
\begin{scope}[shift={(-4,0)}]
\draw[ink,fill=blue!6](-1.8,0)circle(0.35);
\node[]at(-1.80000,-0.80000){源尺度 $a$};
\draw[axis](-1.30000,0.00000)--(2.50000,0.00000);
\node[]at(2.20000,0.35000){观察点 $P$};
\draw[{Latex[length=1.3mm]}-{Latex[length=1.3mm]},ink!65,line width=.5pt](-1.80000,-1.50000)--(2.40000,-1.50000);
\node[lab]at(0.30000,-1.75000){$r$};
\node[]at(0.00000,1.50000){$ka$：源尺寸；$kr$：观察位置};
\node[]at(0.00000,-2.30000){$r\gg a$ 后比较近场与辐射项};
\end{scope}
\node[card]at(4.00000,1.25000){$kL:\text{系统传播尺度};\quad kr:\text{观察距离}$};
\node[card]at(4.00000,-1.05000){$\chi=\omega\varepsilon/\sigma:\text{电流机制}$};
\node[text width=15.2cm,align=center]at(0.00000,-3.20000){源电尺寸记 $ka$；系统传播尺度另用 $kL$，两种长度不可混用。};
\end{tikzpicture}\end{document}
